\section{Security Analysis}\label{sec:security}

This section outlines the security guarantees of our scheme.  
We first analyze the unlinkability and integrity forgery resistance in absence of collusion, and then discuss the effect of potential collusion with STTPs.

%%%%%%%%%%%%%%%%%
% UNLINKABILITY %
%%%%%%%%%%%%%%%%%

\subsection{Unlinkability}\label{sec:unlinkability}

To argue about sender-receiver unlinkability, we assume that an adversary cannot link a sender to its recipient without retracing the entire communication path.  
Thus, the problem reduces to proving unlinkability at the node level, meaning that an adversary cannot determine which outgoing packet corresponds to a given incoming packet.

We formalize unlinkability via a standard indistinguishability game between a challenger $\mathcal{C}$ and an adversary $\mathcal{A}$.
Intuitively, the adversary $\mathcal{A}$ observes packets entering and leaving a node and attempts to determine which outgoing packet corresponds to a given incoming packet.
The scheme provides unlinkability if any probabilistic polynomial-time adversary achieves only a negligible advantage over random guessing, i.e. $\Pr[b_\mathcal{A} = b_\mathcal{C}] - \frac{1}{2}$ is negligible.
% under the hardness of the underlying elliptic curve assumptions.

\medskip
\noindent\textbf{Game Definition:}
\begin{enumerate}[topsep=1pt, itemsep=1pt]
  \item \textbf{Setup.}
  The challenger generates the node's secret keys and public system parameters $(G_1, \dots, G_n)$ as in the setup phase, ensuring all generators $G_j$ are independent and uniformly distributed EC points.
  Generators are said to be independent if no participant knows any scalar relations among them.

  \item \textbf{Oracle.}  
  The adversary is granted access to an oracle, which on input of a nonce $x$ and routing path $\pi = (N_1, N_2, N_3, \Delta)$ outputs a valid header $(\alpha, \beta, \gamma)$ and its updated form $(\alpha', \beta', \gamma')$ after being processed by the node.

  \item \textbf{Challenge.}
  The adversary submits two distinct headers $(\alpha_0, \beta_0, \gamma_0)$ and $(\alpha_1, \beta_1, \gamma_1)$ that were not obtained during oracle queries.  
  The challenger selects a random bit $b_\mathcal{C} \in \{0,1\}$ and processes $(\alpha_{b_\mathcal{C}}, \beta_{b_\mathcal{C}}, \gamma_{b_\mathcal{C}})$ through the node, obtaining $(\alpha'_{b_\mathcal{C}}, \beta'_{b_\mathcal{C}}, \gamma'_{b_\mathcal{C}})$, which is then given to the adversary.

  \item \textbf{Guess.}
  The adversary outputs a bit ${b_\mathcal{A}} \in \{0,1\}$, attempting to guess which input headers corresponds to the given output header.  
  He wins if ${b_\mathcal{A}} = {b_\mathcal{C}}$.
\end{enumerate}

\medskip
\noindent\textbf{Proof Sketch:}
As in the original Sphinx construction, all packets have a fixed size% (header and payload),
, preventing any adversary from correlating packets based on observable size patterns.
In addition, headers are computationally indistinguishable from random bitstrings, as empirically confirmed by the NIST randomness test suite (Section~\ref{sec:NIST}).
Consequently, any successful distinguishing strategy must exploit correlations between header components $(\alpha, \beta, \gamma)$ across layers.
Since the integrity tag $\gamma$ is computed independently at each layer, it cannot be used for cross-layer linkage.

\smallskip
\noindent{\textbf{Case 1: Linking via $\boldsymbol\alpha$.}}  
Network nodes reject packets with duplicate $\alpha$ values, preventing trivial replay-based linkage.
The only relation between $\alpha$ and its processed form $\alpha'$ is the node's shared secret $s$.
Recovering $s$ from $\alpha' = s\alpha$ requires solving an instance of the Elliptic Curve Discrete Logarithm Problem.

\smallskip
\noindent{\textbf{Case 2: Linking via $\boldsymbol\beta$.}}\label{sec:indep_generators}  
Each header chunk $\beta_j$ and its corresponding processed value $\beta'_j$ satisfy $\beta_j - \beta'_j = s G_j$ for all $j \in \{1,\dots,n\}$.
To establish a link, $\mathcal{A}$ must either verify that the same shared secret $s$ is used across all chunks, or determine whether the scalar relationship between two generators $G_j$ and $G_{j'}$ is preserved (i.e. find $z$ such that $G_j = z G_{j'}$).
When generators are independent, it reduces to solving an instance of the ECDLP, either to recover $s$ or to determine $z$.

\smallskip
\noindent\textbf{Summary.}  
Under the ECDLP hardness assumption and assuming independent generators, no probabilistic polynomial-time adversary can achieve a non-negligible advantage in the unlinkability game.
Therefore, we can conclude that the proposed scheme provides unlinkability.

%%%%%%%%%%%%%
% INTEGRITY %
%%%%%%%%%%%%%

\subsection{Integrity Forgery Resistance}\label{sec:integrity_forgery}

Integrity-forgery resistance requires that an adversary cannot modify a header that is accepted by an honest node.  
We formalize this property as a game in which the adversary attempts to forge a new header that passes integrity verification.
The scheme achieves integrity-forgery resistance if any probabilistic polynomial-time adversary has only negligible advantage over randomly guessing the integrity tag.

\medskip
\noindent\textbf{Game Definition:}
\begin{enumerate}[topsep=1pt, itemsep=1pt]
  \item \textbf{Setup.}
  The challenger generates node's secret keys and public system parameters $(G_1, \dots, G_n)$ as in the setup phase, ensuring all generators $G_j$ are independent and uniformly distributed elliptic curve points.

  \item \textbf{Oracle.}
  The adversary is given access to an oracle that outputs valid headers $(\alpha,\beta,\gamma)$.
  The adversary may modify $\beta$ into $\beta'$ and query a second oracle to obtain the corresponding integrity tag $\gamma'$.

  \item \textbf{Challenge.}
  The challenger outputs a fresh header $(\alpha,\beta,\gamma)$ with unused $\alpha$.
  The adversary outputs modified chunks $\beta_{\mathcal{A}}$ and a tag $\gamma_{\mathcal{A}}$.

  \item \textbf{Guess.}
  The adversary wins if $\gamma_{\mathcal{A}} = \gamma_{\mathcal{C}}$, where $\gamma_{\mathcal{C}}$ is the integrity tag computed by the challenger on $\beta_{\mathcal{A}}$.
\end{enumerate}

\medskip
\noindent\textbf{Proof Sketch:}
An adversary is able to modify a header chunk $\beta_j$ in two ways; (i) by replacing it with a new elliptic curve point, or (ii) by adding a point $P_j$.

\smallskip
\noindent\textbf{Case 1: Chunk Replacement.}
Replacing $\beta_j$ with $\beta'_j$ requires updating the integrity tag to $\gamma' = \gamma + h^{(j)}(\sHashed{})\, (\beta'_{j} - \beta_{j})$.
This requires knowledge of the weight $h^{(j)}(\sHashed{})$, which is derived from the hash of the shared secret ($\sHashed{}$). 

\smallskip
\noindent\textbf{Case 2.a: Incremental Modification.}
Adding a point $P_j$ to a single chunk requires updating the tag as $\gamma' = \gamma + h^{(j)}(\sHashed{})\, P_j$, which again requires the weight.

\smallskip
\noindent\textbf{Case 2.b: Balanced Increment.}
An adversary may attempt to modify two chunks $j$ and $j'$ so that the checksum remains unchanged: $\gamma' = \gamma + h^{(j)}(\sHashed{})\, P_j + h^{(j')}(\sHashed{})\, P_{j'}$.
If the scalar relation between two weights is known, i.e. $z$ such that $h^{(j)}(\sHashed{}) = z\,h^{(j')}(\sHashed{})$, it becomes $\gamma' = \gamma + h^{(j')}(\sHashed{})\, (z P_j + P_{j'})$.
Then the adversary can choose any $P_j$ and set $P_{j'} = -z P_j$, leaving $\gamma' = \gamma$ and thus forging integrity.

\smallskip
\noindent\textbf{Summary.}
Integrity forgery requires knowledge of a weight or a scalar relation between weights, which are derived from the hash of a one-time shared secrets.

\medskip
\noindent\textbf{Remarks.}\label{replay_resistance}\label{wrap_resistance}
The original Sphinx provides \emph{replay} and \emph{wrap} resistance \cite{sphinx}.
Replay resistance is guaranteed by dropping previously seen $\alpha$ values.
Wrap resistance prevents transforming a header $(\alpha, \beta, \gamma)$ into $(\alpha', \beta', \gamma')$ that decrypts back to $(\alpha, \beta, \gamma)$.
Both properties require modifying a header that passes integrity verification and thus reduce to integrity forgery resistance.

% \medskip
% \noindent\textbf{Remarks.}
% \label{replay_resistance}\label{wrap_resistance}
% The original Sphinx article~\cite{sphinx} takes \emph{replay resistance}
% and \emph{wrap resistance} into account.
% \emph{Wrap resistance} means an adversary must not be able to transform a header $(\alpha, \beta, \gamma)$ into a new valid header $(\alpha', \beta', \gamma')$ such that, after being decrypted by a node, the packet becomes $(\alpha, \beta, \gamma)$ again.  
% \emph{Replay resistance} is inherited from the original Sphinx construction since the nodes maintain a table of previously seen $\alpha$ values and discard replayed values (with tables reset only when keys are updated).
% Thus, if an adversary want to replay a packet, he would need to alter $\alpha$ (and consequently $\beta$) while still producing a consistent integrity tag $\gamma$. % which directly reduces to the problem of forging integrity.  
% Similarly, \emph{wrap resistance} also requires producing a modified header that passes integrity verification.
% Hence, \emph{replay resistance} and \emph{wrap resistance} reduce to \emph{integrity forgery resistance}. 


%%%%%%%%%%%%%
% COLLUSION %
%%%%%%%%%%%%%

\subsection{Collusion Analysis}\label{sec:collusion}

The preceding security arguments assume no collusion. 
Each STTP learns additive shares of destinations and per-hop shared secrets, which reveal no information unless all shares are combined.
In addition, STTPs receive the values $\sigma_i = h(s_i)$ for all hops in the path, from which integrity weights are derived. 
However, the blinding offset $s_iG$ prevents associating these values with a specific packet. 
Consequently, collusion among a subset of STTPs does not compromise unlinkability or integrity.

The main risk arises when an STTP colludes with a mixnode. 
A corrupted mixnode can derive its shared secret $s_i$ from $\alpha$ and compute the corresponding hash $\sigma_i$, which can then be matched against the set of values observed by the STTP.
If a match is found, the adversary can associate the packet with the set of values $\{\sigma_j\}$ used along its path. 
This enables integrity forgery and weakens unlinkability by allowing probabilistic packet tracing. 
To mitigate this threat, the output space of $\sigma$ can be reduced, for example by truncating the hash output, creating multiple candidate matches and preventing unique packet-weight associations.
Although this increases the probability of random guessing, the blinding offset $s_iG$ prevents efficient guess verification. 
Therefore, security is preserved as long as at least one STTP remains honest, while the impact of STTP-mixnode collusion can be limited through an appropriate choice of integrity-weight space, which balances collusion resilience against random-guessing resistance.
Determining the optimal trade-off between collusion resilience and guessing resistance is left for future work.
